\begin{abstract}
An agent partway through a long task cannot say how much more it needs. Asked mid-task for an interval on its remaining tokens, Claude~Opus~4.7 contains the truth 51\%, 23\%, and 35\% of the time on Sokoban, Search-R1, and SWE-bench, against a 90\% target. Budget Guard supplies that number from outside the agent. It reads the per-turn usage counts every provider API already reports, fits a small model to them, and wraps that model in conformal calibration, so that its interval on the remaining requirement contains the truth as often as promised. Nothing in the agent is trained or changed. Budget awareness turns out to be two questions, how much more and whether to stop, and we build a separately calibrated instrument for each. The intervals reach 0.90--0.96 coverage offline and 0.87--0.92 online with no unbounded intervals, hold at least 0.85 in 18 of 19 model--environment runs on one split each, and are sharper than a feature-free baseline; that much is guaranteed in distribution, and the experiments map what happens outside it. A guard moved from Sokoban to SWE-bench without recalibration falls to 0.10 coverage; recalibrating from finished runs recovers 0.80, and no bounded method reaches the target after a model swap. When three budgets bind at once, per-resource calibration covers the joint event 80\% of the time; a Bonferroni correction restores 0.95 in distribution. The stopping instrument, a feasibility gate judged on money in Warehouse replay, avoids 86.3\% of the spend on runs that could not finish, stops no run that would have finished, and cuts cost per success from \$1.61M to \$0.87M in simulated dollars, which no hard cap approaches. Feeding the intervals into the stop rule only makes it more conservative: the intervals measure, the gate decides. All stopping results are counterfactual replays of logged runs.
\end{abstract}

\raggedbottom
\section{Introduction}
\label{sec:intro}

LLM agents now take on tasks that run for many steps, many tool calls, and many interactions with an environment~\cite{kwa2025measuring}, and the resources they consume grow with them~\cite{chen2026token}. A single attempt can burn thousands of tokens, real time, and real money before it ends~\cite{bai2026how}, much of it on reasoning that changes nothing~\cite{sui2025stop,wang2025harnessing,yue2025dont}; when the attempt fails, nearly all of it was wasted~\cite{bai2026how}. Keeping that spending under control without giving up task performance is one of the practical obstacles to running agents at scale~\cite{chen2024frugalgpt,zellinger2025cost,davidov2026knowing}.

The obvious fix is to ask the agent, and it does not know. BAGEN~\cite{lin2026bagen} asked agents mid-task for an interval on how many more tokens they would need. For Claude~Opus~4.7 those intervals contain the true number 51\%, 23\%, and 35\% of the time on Sokoban, Search-R1, and SWE-bench, against a 90\% target. Figure~\ref{fig:case_study} shows what that looks like on one run: the agent's interval looks precise and sits above the truth from the second turn on, reserving 400 to 900 tokens for a run that needed 267.

\begin{figure}[H]
\centering
\includegraphics[width=0.72\linewidth]{figures/fig_case_study.png}
\caption{What Budget Guard does on one real run: the first held-out Sokoban rollout in our seed-0 test stream (Claude~Opus~4.7, 2{,}500-token budget), chosen by that rule rather than by hand. Grey bars are the agent's own interval on the tokens it still needs, green bars the guard's, black dots the truth, and the dashed line the budget still available. The agent's interval misses from turn~2; the guard's contains the truth at every turn and stays under the remaining budget, and the guard's separate stop-or-continue gate (scores under each turn) says continue throughout. The run succeeded inside budget. Table~\ref{tab:case_study} in Appendix~\ref{app:breadth} gives the numbers.}
\label{fig:case_study}
\end{figure}

Budget Guard supplies the missing number from outside the agent. It is a wrapper: it reads the per-turn usage that provider APIs already report, fits a small gradient-boosted model to it, and wraps that model in a conformal calibration layer, so that its interval on the remaining requirement contains the truth as often as promised. Nothing in the agent is trained or changed, and because the guard reads only usage counts it works for closed models whose internals are unavailable. In Figure~\ref{fig:case_study} the green intervals are wider than the agent's and right every time; spend so far plus their upper bound stays under the budget, so a deployer could have certified from the first turn that the run fit. That is the first half of budget awareness: how much more.

The second half is whether to stop, and it is a different question with a different instrument. A feasibility gate on the same features scores each turn and aborts a run when the score drops below a threshold calibrated so that at most a chosen fraction $\delta$ of runs that would have finished are stopped. We keep the two apart throughout, with separate calibration targets, separate metrics, and separate claims, because the results say we must: the intervals are a reliable measurement, and using them as a stopping signal only makes the gate more timid.

Two facts about the label shape everything. The resources a run still needs are known only when it ends, so online methods may read a label only at rollout boundaries; an earlier revision of this evaluation read them per prefix, and correcting that lowered our headline cross-environment recovery from 0.90 to 0.80. And the label does not exist at all for a run that can no longer finish, which is exactly the run a guard exists to catch. Half of our Warehouse prefixes are in that state.

Our findings, in order. (1)~In distribution the intervals do what the construction promises, 0.90--0.96 coverage offline and 0.87--0.92 online, with no unbounded intervals and sharper bounds than a feature-free baseline on two of three environments; the agent's own estimate adds nothing as a feature (\S\ref{sec:within}). (2)~A frozen guard collapses under an environment change, from 0.92 to 0.10, and rollout-safe recalibration recovers 0.71 and 0.80 with bounded intervals; the methods that recover more do so with unbounded intervals on 60--88\% of the stream, and nothing bounded reaches the target after a model swap (\S\ref{sec:shift}). (3)~When several budgets bind at once, per-resource calibration under-covers the joint event, and a Bonferroni correction restores it in distribution but not under swap (\S\ref{sec:joint}). (4)~The gate avoids 86.3\% of doomed-run spend in Warehouse replay with no false aborts, which no hard cap achieves; on token tasks it saves far less, and adding the intervals to the rule makes it worse (\S\ref{sec:econ}). (5)~Runs that cannot finish carry information that dropping them throws away (\S\ref{sec:censored}). Table~\ref{tab:guarantees} says which of these are guarantees and which are observations, and every number regenerates from one command over a pinned dataset.

\section{Related work}
\label{sec:related}

BAGEN~\cite{lin2026bagen} defines budget awareness as an agent's ability to estimate mid-execution what it still needs, and asks the agent directly for an interval on remaining tokens and a judgment of whether finishing is still feasible, over internal resources such as tokens and external ones such as time, money, and warehouse capacity, on Sokoban, Search-R1, SWE-bench, and Warehouse. Its central finding is a gap between capability and self-knowledge: strong agents solve hard tasks while remaining poorly calibrated about what finishing will cost. Stopping early on the agent's own signals saves roughly 28--64\% of tokens on failed trajectories at a cost of roughly 1.6--4.2 percentage points of task success, and interval coverage reaches only 47\% even after supervised fine-tuning and reinforcement learning~\cite{lin2026bagen}. The agent carries some signal about its future; its raw estimate is not reliable enough to act on. Conformal prediction~\cite{angelopoulos2023gentle,lei2018distribution} is the standard tool for turning an unreliable predictor into an interval with a stated coverage rate, and that is what we use it for.

Doomed From the Start~\cite{ruan2026doomed} reads a model's internal activations to decide that a trajectory is doomed and stops it early. Like us, and unlike BAGEN, it uses a distribution-free calibrated gate with no fine-tuning, meeting recall targets of 90--97\% on TextCraft and saving 47.1\% (Qwen-2.5-7B) and 37.2\% (Llama-3.2-3B) of inference compute at the 90\% target. Because it needs internal states it applies only to open-weight models; Budget Guard needs only usage counts. MiCP~\cite{zhou2026adaptive} applies conformal prediction to a different object, the agent's answer, allocating an error budget across the turns of a multi-turn reasoning process so the agent can stop once it has enough evidence. We calibrate coverage over the remaining resource instead, sit outside the reasoning loop, and keep the abort gate separate from the intervals. Section~\ref{sec:method} specifies the construction; Appendix~\ref{app:background} defines the terms used to read the results.

\FloatBarrier
\section{Method}
\label{sec:method}

\paragraph{The interval.}
Take any predictor, measure how badly it missed on data it never trained on, and widen its output by a high quantile of those misses. We fit two quantile regressors, at $\alpha/2$ and $1{-}\alpha/2$, giving a band $[\hat f_{\mathrm{lo}}(x), \hat f_{\mathrm{hi}}(x)]$; on each of $n$ calibration points we record how far the truth fell outside it, $s=\max(\hat f_{\mathrm{lo}}-y,\ y-\hat f_{\mathrm{hi}})$, negative when the point sat inside; and with $\hat q$ the $\lceil (n+1)(1{-}\alpha)\rceil/n$ empirical quantile of those scores, the interval $[\hat f_{\mathrm{lo}}(x)-\hat q,\ \hat f_{\mathrm{hi}}(x)+\hat q]$ contains $y$ with probability at least $1{-}\alpha$~\cite{romano2019cqr}. We use $\alpha{=}0.10$ and clip the lower endpoint at zero. Two things about this guarantee matter for everything that follows. It is a hit rate over runs, not a property of any one interval: across exchangeable runs the interval contains the truth at least 90\% of the time, and it says nothing about which 10\% it misses. And it asks nothing of the data except that the calibration runs and the deployment runs be exchangeable. Everything difficult in this paper comes from the fact that they often are not. Table~\ref{tab:guarantees} records, for each component, exactly what we claim and under what assumption.

\begin{figure}[t]
\centering
\includegraphics[width=\linewidth]{figures/pipeline.pdf}
\caption{Budget Guard wraps an unmodified agent in an external calibration layer. Stages~1--4 share prefix features; stage~3 returns a calibrated remaining-budget interval (\S\S\ref{sec:within}--\ref{sec:joint}), and a separately calibrated feasibility score drives continue-or-abort (\S\ref{sec:econ}). Warehouse abort economics use the feasibility gate alone and do not read the intervals. Labels become available only when a rollout finishes, so online recalibration never uses a still-running run.}
\label{fig:pipeline}
\end{figure}

\paragraph{What the guard sees.}
At turn $k$ the guard sees only the past (Figure~\ref{fig:pipeline}): how far along the run is (the turn index); how much it has consumed (spend so far, and spend per turn, the burn rate); how it has been consuming recently (the last turn, the mean of the last three, and the mean, standard deviation, minimum, maximum, and trend of per-turn spend); and the agent's own estimate when it gives one. Nothing that depends on the future (final totals, success, feasibility) enters the feature matrix; we assert this in code rather than trusting ourselves to have left it out. No weights, no logits, no environment internals: only the usage accounting that provider APIs already return, which is what makes the guard a wrapper rather than a modification, and a verifier rather than a believer. It never has to trust the agent's own estimate, because it wraps that estimate in a guarantee that holds whether the estimate is any good or not. The predictor behind the band is gradient-boosted quantile regression~\cite{ke2017lightgbm}, so the width varies with the input. We split calibration and test by rollout, never by prefix: two prefixes of one run share its final total, and splitting them apart would leak each test answer into calibration.

\paragraph{Recalibrating as the stream moves.}
Frozen split-CQR is valid while calibration and deployment are exchangeable. Swap the underlying model, or point the guard at a new task type, and the resemblance is gone; so is the guarantee. Our response is to let the deployment stream recalibrate the guard from the runs that have already finished, in three ways. ACI~\cite{gibbs2021aci} nudges its working confidence level up when it has been missing too often and down when it has been too safe. DtACI~\cite{gibbs2024dtaci} runs several such nudges at different speeds and leans on whichever is currently working. Windowed-CQR keeps a sliding window of the 100 most recent scores and reads $\hat q$ straight off them; the idea is borrowed from SPCI~\cite{xu2023spci}, though we take the plain empirical quantile rather than fitting their model of the residual series, and the honest statement of what it provides is a coverage gap that grows with how far the window has drifted from the stream~\cite{barber2023beyond}. All three see a label only when its rollout ends, and hold their state fixed within a run. Two reporting rules follow. An interval with no upper bound covers everything and tells you nothing, so the fraction of unbounded intervals sits beside every coverage number. And because one margin is set by the many easy late-turn prefixes, a good average can hide a hole at turn one, where an abort decision is worth the most; we therefore also calibrate a separate $\hat q$ per turn-depth bucket, $k{=}1$, $k{=}2$, $k{\ge}3$ (Mondrian~\cite{vovk2012conditional}).

\paragraph{Several budgets, and runs that cannot finish.}
A Warehouse run can fail three ways, out of time, capacity, or money, and it survives only if all three estimates hold. Covering each 90\% of the time does not cover all three 90\% of the time, because the misses accumulate; calibrating each at miscoverage $\alpha/3$ (Bonferroni) restores the joint guarantee by the union bound at the price of wider margins, and is our primary scheme. A tighter alternative shares one threshold across resources by ranking their errors jointly~\cite{messoudi2021copula}, but that rank is chosen on the same errors it then thresholds, so it carries no guarantee; we report it as an observation, and let both move under a model swap from a rolling window of finished runs. When a run is doomed we never learn $y$, only that it exceeded what was left, and dropping such runs removes exactly the hard cases. We calibrate one-sided upper bounds five ways (censored scores as $+\infty$, always valid and vacuous once censoring exceeds $\alpha$; the known lower bound substituted; censored runs dropped, the naive comparison; reweighting by censoring probability, IPCW~\cite{robins1992recovery,gui2024conformalized}; and charging a censored run only against the resource that stopped it), and a classifier on the same features can decline to issue an interval at all. That changes what we can promise: coverage among the prefixes the gate lets through, not across all prefixes. It is weaker than a marginal guarantee, but it is what the construction supports, so it is what we claim.

\paragraph{The gate.}
The same feasibility score decides when to quit: abort at the first turn where it drops below a threshold $\tau$. Since a run is aborted the moment its lowest score crosses $\tau$, we set $\tau$ from the lowest scores of the feasible calibration runs, taking their $\delta$-quantile with the finite-sample correction above, which bounds the probability of aborting a run that would have succeeded by $\delta$, a per-run guarantee in the sense of conformal risk control~\cite{angelopoulos2024crc}. Aborting a doomed run recovers what it had left to spend; aborting a good one throws away a success. Both land in the same currency, so the gate is judged on Cost-of-Pass~\cite{erol2026costofpass}, total spend per successful run, rather than on coverage. The gate reads the feasibility score and nothing else; a variant that also aborts when spend so far plus the interval's calibrated lower bound exceeds the budget is evaluated on the token environments (Table~\ref{tab:abort}).

\section{Experimental setup}
\label{sec:setup}

\paragraph{Data and environments.}
All experiments use logged rollouts from BAGEN~\cite{lin2026bagen} at a pinned dataset revision; Table~\ref{tab:envs} summarises the four environments. Sokoban, Search-R1, and SWE-bench are single-resource token tasks. Sokoban and Search-R1 use uniform context budgets of 2{,}500 and 3{,}500 tokens; on SWE-bench the official budget varies by run (6{,}772, 7{,}060, 8{,}265, 10{,}472, and 21{,}711 tokens across the five models) and we read each run's value from its own record. The label is the number of tokens still needed to complete the task under BAGEN's history-excluded per-turn token accounting. Warehouse is a multi-week management simulation with three jointly binding budgets (time in weeks, capacity in item-weeks, cost in simulated USD), with per-prefix ground truth for all three from the official records; half of its prefixes (640 of 1{,}280 per model) belong to runs that can no longer finish, which makes it our censoring test.

\begin{table}[t]
\centering
\caption{The four evaluation environments; only Warehouse binds several resources at once, and only Warehouse censors labels. Token budgets are the BAGEN context budgets in tokens; Warehouse budgets are set per task. In Warehouse, half of the prefixes (640 of 1{,}280 per model) belong to runs that can no longer finish and drive the censored-calibration evaluation.}
\label{tab:envs}
\small
\begin{tabular}{@{}llllc@{}}
\toprule
Environment & Task & Budgeted resource(s) & Budget & \# Models \\
\midrule
Sokoban & box-pushing puzzle & tokens & 2{,}500 tokens & 7 \\
Search-R1 & multi-turn retrieval QA & tokens & 3{,}500 tokens & 7 \\
SWE-bench & repository-level bug fixing & tokens & 6{,}772--21{,}711 (per run) & 5 \\
Warehouse & multi-week management sim. & time, capacity, cost & per task, joint & 5 \\
\bottomrule
\end{tabular}
\end{table}

\paragraph{Protocol, baselines, and models.}
Rollouts are split 50/25/25 into train, calibration, and test, always by whole rollout. Every primary experiment is repeated over five random splits (seeds 0--4) and reported as the mean; the 19-run breadth count is one split per run, and we say so wherever it appears. Online methods consume the test stream in order with labels revealed only at rollout boundaries. We evaluate three shift scenarios: cross-environment (calibrate on Sokoban; deploy on Search-R1 and SWE-bench), model swap (calibrate on Claude~Opus~4.7; deploy on four other agents' SWE-bench streams), and an external-resource model swap on Warehouse. Seven comparators run on the same splits and seeds as the guard: the agent's own interval; a feature-free constant-quantile interval, which measures what the learned component adds; split conformal on absolute residuals; the agent's own midpoint conformalized; the guard with the agent's estimate added to its features; frozen split-CQR under shift; and hard caps for stopping. Coverage is always reported against the 0.90 nominal level and beside the fraction of unbounded intervals. We evaluate seven agents (Claude~Opus~4.7/4.6, Claude~Sonnet~4.6, GPT-5.2-instant, GPT-5.2-Codex, Gemini~3.1~Pro, and Qwen3-235B), which crossed with environments gives 19 token-budget model$\times$environment runs (7 Sokoban, 7 Search-R1, 5 SWE-bench) of 52--128 rollouts and 85--3{,}715 prefix samples each. Appendix~\ref{app:protocol} gives the full protocol, baseline definitions, metric formulas, and implementation details.

\section{Results}
\label{sec:results}
\raggedbottom

\subsection{Within-run interval calibration}
\label{sec:within}

\begin{figure}[t]
\centering
\includegraphics[width=0.92\linewidth]{figures/fig_coverage.png}
\caption{Conformal calibration closes almost all of the gap between agent self-reports and the 0.90 target, and windowed-CQR does so while maintaining bounded intervals throughout. Marginal coverage of the true remaining budget on Sokoban, Search-R1, and SWE-bench (Claude~Opus~4.7 rollouts), $\alpha{=}0.10$, averaged over five rollout-level split seeds; error bars are the standard deviation across those seeds. Hatched bars emit unbounded intervals on the annotated fraction of the stream, which raises coverage at the cost of informativeness. Agent self-report is a single value read from the logged records rather than a seed mean. Coverage, normalised width, unbounded fraction, and test-set sizes for all methods are in Table~\ref{tab:coverage}.}
\label{fig:coverage}
\end{figure}

\begin{figure}[t]
\centering
\includegraphics[width=\linewidth]{figures/fig_allruns.png}
\caption{Every token-budget run, one split each (seed~0), no averaging: coverage of the agent's own interval and of split-CQR, windowed-CQR, and DtACI for each of the 19 model--environment runs, against the 0.90 target and the 0.85 line of the breadth claim. Windowed-CQR clears 0.85 on 18 of 19 runs and 0.90 on 11; the agents' own intervals clear neither anywhere. GPT-5.2-Codex on Sokoban carries no agent interval in its records. Table~\ref{tab:allruns} lists the values.}
\label{fig:allruns}
\end{figure}

Figure~\ref{fig:coverage} is the in-distribution picture on Claude~Opus~4.7. The agent's own intervals cover 0.51, 0.23, and 0.35 of prefixes on Sokoban, Search-R1, and SWE-bench against a 0.90 target; offline split-CQR reaches 0.92, 0.90, and 0.96, and online windowed-CQR 0.92, 0.92, and 0.87, with no unbounded intervals anywhere. That split-CQR lands near 0.90 when calibration and test are exchangeable is what the construction guarantees; we report it as a check, not a result, and the findings are everything the guarantee does not cover. The intervals are bounded and sharp: windowed-CQR is 22\% narrower than the feature-free constant-quantile baseline on Sokoban and 55\% narrower on SWE-bench, though 23\% wider on Search-R1, where constant-quantile is the sharpest bounded method (Table~\ref{tab:coverage}). ACI-CQR and DtACI sometimes post higher coverage, partly by emitting unbounded intervals on up to 16\% of the stream. The agent's own estimate adds nothing as a feature: guard+self-est is 0.91 against split-CQR's 0.92 on Sokoban, identical on SWE-bench (0.96), and worse on Search-R1 (0.85 against 0.90). Figure~\ref{fig:allruns} shows every run rather than one agent: across all 19 logged model--environment runs, windowed-CQR reaches at least 0.85 coverage in 18 of 19 and at least 0.90 in 11 of 19, against 14 and 8 for DtACI, on one split per run, and the agents' own intervals clear neither on any run (Table~\ref{tab:allruns}). The run of Figure~\ref{fig:case_study} is typical rather than selected: over the 30 test rollouts of that split, split-CQR covers more prefixes than the self-report on 22, ties on 6, and trails on 2 (windowed-CQR: 24, 5, and 1).

\enlargethispage{2\baselineskip}
\subsection{Robustness under deployment shift}
\label{sec:shift}

\begin{figure}[t]
\centering
\includegraphics[width=\linewidth]{figures/fig_transfer.png}
\caption{An environment change collapses frozen calibration and online recalibration recovers much of it; a model swap leaves every method below nominal coverage. Marginal coverage on the target stream. (a)~A guard calibrated on Sokoban and deployed unchanged on Search-R1 and SWE-bench: frozen split-CQR falls to 0.51 and 0.10, and windowed-CQR recovers to 0.71 and 0.80. (b)~A guard calibrated on Claude~Opus~4.7 SWE-bench rollouts and deployed on four other agents: here frozen and online calibration are comparable on average (0.74 against 0.73) and neither reaches the target, though windowed-CQR varies far less across seeds. Hatched bars emit unbounded intervals on the annotated fraction of the stream (60--88\% for the ACI-based methods in~(a), and 22--61\% (ACI-CQR) and 4--7\% (DtACI) in~(b)).}
\label{fig:transfer}
\end{figure}

This is where the guarantee stops and the experiment starts. Figure~\ref{fig:transfer}(a) takes a guard calibrated on Sokoban and points it, unchanged, at Search-R1 and SWE-bench: frozen split-CQR falls to 0.51 and 0.10. Rollout-safe windowed-CQR, recalibrating only from runs that have finished, recovers 0.71 and 0.80 with no unbounded intervals. ACI and DtACI post higher numbers but need unbounded intervals on 60--88\% of predictions to get there, which is recovery by surrender rather than by adaptation. Figure~\ref{fig:transfer}(b) swaps the agent instead of the task: a guard calibrated on Claude~Opus~4.7 is deployed on Sonnet, Gemini, GPT, and Qwen SWE-bench streams. Windowed-CQR reaches 0.66--0.89, below target everywhere, but its seed-to-seed standard deviation is an order of magnitude smaller than frozen calibration's (0.006--0.011 against 0.122--0.224), so its failures are at least predictable. Online adaptation helps under shift and does not solve it: no bounded-interval method recovers the nominal level under a model swap, and any coverage number has to be read next to how many of its intervals were infinite.

\subsection{Conditional and multi-resource calibration}
\label{sec:joint}

A marginal number can hide a hole in the slice that matters. On Sokoban, split-CQR reaches 0.92 marginal coverage but only 0.86 at the first turn, exactly where an abort decision has the most leverage; Mondrian calibration by turn depth raises first-turn coverage to 0.93, is neutral to slightly negative on Search-R1, and neutral to slightly positive on SWE-bench, where the first-turn gap is already small (Table~\ref{tab:conditional}). What it buys is conditional validity, not a uniform improvement.

\begin{table}[H]
\centering
\caption{Independent margins under-cover the joint event, Bonferroni fixes it in distribution, and nothing with a guarantee survives a model swap. Warehouse coverage against a nominal 0.90, mean over five rollout-level split seeds; above 0.90, coverage is bought with wider margins. Per-resource columns and the in-distribution joint column are Claude~Opus~4.7; the final column is the joint event averaged over four swap targets. Bold marks the highest joint coverage in each column. ``--'' marks a scheme not run in that setting.}
\label{tab:joint}
\small
\begin{tabular}{@{}lccccc@{}}
\toprule
& \multicolumn{3}{c}{per-resource} & \multicolumn{2}{c}{joint} \\
\cmidrule(lr){2-4}\cmidrule(lr){5-6}
Calibration scheme & time & capacity & cost & in-dist. & model swap \\
\midrule
agent self-report & 0.87 & 0.67 & 0.71 & 0.42 & 0.34 \\
independent margins ($\alpha$) & 0.99 & 0.89 & 0.92 & 0.80 & -- \\
Bonferroni ($\alpha/3$) & 1.00 & 0.98 & 0.97 & \textbf{0.95} & 0.71 \\
empirical copula & 0.99 & 0.94 & 0.96 & 0.89 & 0.62 \\
adaptive copula (online) & -- & -- & -- & -- & \textbf{0.87} \\
\bottomrule
\end{tabular}
\end{table}

Table~\ref{tab:joint} is the multi-resource case. In Warehouse the agent's self-reports cover the joint event across time, capacity, and cost 0.42 of the time. Calibrating each resource independently at $\alpha$ lifts that only to 0.80: independent margins do not compose into joint coverage. Bonferroni at $\alpha/3$ restores 0.95, and the empirical copula reaches 0.89 with narrower margins but no guarantee. Under an external-resource model swap, frozen Bonferroni falls to 0.71, below its nominal level and below the 0.87 of the un-guaranteed adaptive copula, which recalibrates from finished rollouts. We report both and rest the claim on the one with the guarantee.

\subsection{Economic impact}
\label{sec:econ}

\begin{figure}[t]
\centering
\includegraphics[width=\linewidth]{figures/fig_economics.png}
\caption{Warehouse economic impact from the feasibility gate alone, without remaining-budget intervals. At $\delta{=}0.10$ the classifier-only rule avoids 86.3\% of doomed-run spend with no observed false aborts, which no hard cap achieves at any setting. Simulated USD, mean over five seeds. (a)~Every hard cap that saves more than 0.2\% of doomed spend also stops runs that would have succeeded; the gate occupies the favourable region of the savings--false-abort trade-off. (b)~Cost-of-Pass against the fraction of successful runs retained: the gate roughly halves cost per pass while retaining every observed success, whereas no cap improves on the unstopped baseline. Both panels are counterfactual replays of logged rollouts, not live deployments.}
\label{fig:econ}
\end{figure}

The gate is judged on money, not coverage. Figure~\ref{fig:econ} replays Warehouse with the classifier-only gate at $\delta{=}0.10$: it avoids 86.3\% of doomed-run spend with 0\% observed false aborts. Hard spending caps cannot: the tightest dollar cap that preserves every successful run saves 0.2\% of doomed spend, dollar caps reach meaningful savings only at false-abort rates of 19--55\%, and the only week cap that saves anything (54.5\%) aborts every feasible run. The gate cuts Cost-of-Pass from \$1.61M to \$0.87M, a reduction of 46\%, while retaining every observed success; no cap improves on the unstopped baseline by more than a rounding error, and the tightest raises cost per pass to \$3.19M by discarding passes.

\begin{table}[H]
\centering
\caption{Abort performance at $\delta{=}0.10$ (mean over five seeds). Warehouse uses the classifier-only gate (native USD). Token abort is evaluated on Sokoban and SWE-bench only; both also evaluate a classifier+interval variant. Combining intervals reduces doomed savings at this operating point.}
\label{tab:abort}
\small
\begin{tabular}{@{}llccc@{}}
\toprule
Environment & Rule & False abort & Doomed spend saved & Uses intervals? \\
\midrule
Warehouse & feasibility only & 0.0\% & 86.3\% & No \\
Sokoban & feasibility only & 5.8\% & 7.1\% & No \\
Sokoban & feasibility+interval & 5.8\% & 1.9\% & Yes \\
SWE-bench & feasibility only & 0.0\% & 3.1\% & No \\
SWE-bench & feasibility+interval & 0.0\% & 0.0\% & Yes \\
\bottomrule
\end{tabular}
\end{table}

Table~\ref{tab:abort} puts the same gate on the token tasks, and the picture changes. At the same $\delta{=}0.10$, Sokoban and SWE-bench save only 7.1\% and 3.1\% of doomed token spend, with false-abort rates of 5.8\% and 0\%. Adding the interval trigger to the rule makes it more conservative at fixed $\delta$ (Sokoban savings fall to 1.9\%, SWE-bench to 0\%) and recovers none of the Warehouse-scale savings. This is the result behind keeping the two components apart. The intervals measure how much a run still needs, and \S\ref{sec:within} shows they measure it reliably; they are not a stopping signal, and used as one they only make the gate timid. What they are for is planning and monitoring: a deployer reads spend so far plus the upper bound against the budget to see whether a run still fits, as in Figure~\ref{fig:case_study}, or sizes the quota for a batch of runs before launching it. The gate is for stopping. The 86.3\% is therefore a Warehouse result for the gate alone, not a general Budget Guard savings rate. Two further cautions: the dollars are simulated Warehouse dollars, not comparable to BAGEN's 28--64\% token savings; and an observed 0\% false-abort rate on a small test set is consistent with the $\delta{=}0.10$ bound, not proof that the true rate is zero.

\subsection{Censored calibration}
\label{sec:censored}

Roughly half of Warehouse prefixes are censored: the run can no longer finish, so what it would have needed is never observed. Treating those runs as missing data is the wrong call. Dropping them yields only 0.50 coverage on the censored prefixes where cost is the binding resource, which are exactly the expensive failures. Treating censored scores as $+\infty$ is valid but vacuous at this censoring rate, and the informative treatments trade validity for sharpness in ways Table~\ref{tab:censored} makes explicit. The other lever is to abstain: feasibility-gated selective intervals hold coverage at 0.92--0.95 for capacity and cost, and 1.00 for time, on the prefixes they keep, at the cost of keeping between 51\% and 89\% of them (Table~\ref{tab:gated}). That gate is swept over its own prefix-level tolerance $\delta_{\mathrm{gate}}$, a different quantity from the rollout-level $\delta$ of \S\ref{sec:econ}. One caveat governs these numbers: the recovered labels used to score censored prefixes are exact only for time, and carry a week-proportional interpolation with 12--27\% median relative error for capacity and cost, so those two are coverage against an approximate label rather than ground truth.

\flushbottom
\section{Discussion and limitations}
\label{sec:discussion}

The paper shows that an agent's own remaining-budget estimate can be replaced, from outside, by an interval that is right as often as promised in distribution and degrades visibly rather than silently under shift; that no bounded method reaches the target under a model swap; that joint coverage over several budgets needs a joint correction and loses its guarantee under swap; and that a separately calibrated gate, not the intervals, is what saves money, and only where doomed runs are expensive and identifiable early. It does not show a general savings rate: the 86.3\% is Warehouse, gate-only, and falls to 7.1\% and 3.1\% on token tasks at the same $\delta$ (Table~\ref{tab:abort}). The joint results rest on one multi-resource environment, which cannot support a claim that they generalise. The in-distribution coverage numbers should be read as a check on the construction; the evidence for the method is in the shift, sharpness, joint, censoring, and economic results.

The breadth count is a single split: on seed~0, windowed-CQR reaches at least 0.85 coverage in 18 of 19 runs. Search-R1 leaves about 24 test prefixes per seed, which is why its standard deviations reach 0.10--0.17. Coverage alone does not distinguish a useful interval from a trivially wide one, so width and unbounded fraction are reported throughout; Figure~\ref{fig:case_study} also shows the price of calibration, intervals roughly twice as wide as the agent's, which costs nothing when the question is whether the run still fits and matters when a tight point estimate is needed. Censored capacity and cost results are measured against an interpolated label. Finally, savings and Cost-of-Pass are counterfactual replays: they assume adding the guard would not change agent behaviour and that the logged sample is representative.

\paragraph{Broader impact.}
The intended effect is to reduce wasted compute, energy, and money on agent runs that cannot succeed; the Warehouse gate at $\delta{=}0.10$ illustrates that potential in one environment (86.3\% of doomed-run spend avoided in replay, without using remaining-budget intervals). The corresponding risk is that a stopping threshold set too aggressively terminates runs that would have finished; $\delta$ calibration exists to bound that risk at a user-chosen level, and the counterfactual caveat above means a deployer should not treat the reported savings as a guaranteed live return when choosing $\delta$, nor assume token environments will see Warehouse-scale savings.

\section{Conclusion}
\label{sec:conclusion}

Budget awareness can be supplied from outside the agent. A training-free layer over the usage counts an API already reports turns them into remaining-budget intervals with 0.90--0.96 coverage offline and 0.87--0.92 online on Claude~Opus~4.7 against a 0.90 target, and at least 0.85 in 18 of 19 model$\times$environment runs on one split each, where the agent's own intervals manage 0.23--0.51. That guarantee assumes exchangeability, and the experiments map what happens without it: online recalibration recovers much of an environment shift, no bounded method reaches nominal under a model swap, and joint calibration works in distribution and not under swap. A separately calibrated feasibility gate avoids 86.3\% of doomed-run spend in Warehouse replay at $\delta{=}0.10$, and far less on token tasks. Budget awareness is a measurement and a decision, and it works when the two are calibrated, evaluated, and claimed separately.


